\section{الفصل~\ref{sec:synthesis}. الآلة كاملة}

لا يعرض فصل الخلاصة تجارب جديدة: كل سطر فيه يستند إلى الفصل الذي نوقش فيه الادعاء، وإلى المصادر نفسها؛ الراسخ هو ناقل \nt{GLUT} والتحكم في التدفق عبر \nt{GABA}، أما أدوار قنوات البث وعملها المشترك ففرضية في معظمها.

{\footnotesize
\setlength{\tabcolsep}{3pt}\renewcommand{\arraystretch}{1.15}
\begin{longtable}{>{\raggedright}p{0.5cm}>{\raggedright}p{4.0cm}>{\raggedright}p{5.6cm}>{\raggedright}p{2.3cm}>{\raggedright\arraybackslash}p{1.7cm}}
\toprule
م & القضية & التجربة وما قيس & المصدر & الحالة \\
\midrule\endfirsthead
\toprule
م & القضية & التجربة وما قيس & المصدر & الحالة \\
\midrule\endhead
1 & لـ$8.6\cdot10^{10}$ عصبون أربعة أنواع فقط من إشارات التحكم~\eqref{eq:machine} & عدّ نوى الخلايا في متجانس الدماغ (المجزِّئ المتساوي الخواص): لدى الرجل البالغ $86.1\pm8.1$ مليار عصبون & \cite{Azevedo2009} & مُقاس \\
2 & القضية~\ref{prop:architecture}، الطبقتان الأوليان: تسير البيانات على ناقل العناوين \nt{GLUT}، ويحدد المرسلُ علامةَ الوصلة، وتوجّه عصبونات \nt{GABA} التدفق وتطبّعه (الفصلان~\ref{sec:neurontypes} و\ref{sec:gaba}) & ناقل عصبي رئيسي واحد لكل عصبون (مراجعة للقياسات)؛ التثبيط الأمامي يضيّق النافذة التي يجمع فيها العصبون الهرمي مداخله؛ والقسمة على النشاط الكلي مقيسة في مناطق كثيرة & \cite{Strata1999,Pouille2001,Carandini2012} & مُقاس \\
3 & العناصر غير موثوقة: يفقد المشبك معظم النبضات (الجدول~\ref{tab:computer}، الفصل~\ref{sec:code}) & شرائح الحُصين، تنبيه أدنى: أكثر من نصف النبضات لا يستدعي استجابة في \foreignlanguage{english}{CA1}؛ استُبعد إخفاق العتبة والتوصيل على المحور، والسبب هو التحرير الاحتمالي للناقل العصبي & \cite{AllenStevens1994} & مُقاس \\
4 & يعمل الدماغ على $\sim20$ واطًا، بمعدل نبضة في الثانية تقريبًا لكل عصبون (الفصل~\ref{sec:code})؛ ولم يبنِ المهندسون بعدُ آلة كهذه & التقدير~\eqref{eq:energy} من التكاليف المقيسة: نحو $10^9$ جزيء \foreignlanguage{english}{ATP} لكل نبضة مع عمل المشابك (قشرة القوارض)؛ والتقدير المفصّل للقشرة يعطي ترددًا أدنى. حالة التقنية بحسب مراجعة & \cite{Attwell2001,Lennie2003,MehonicKenyon2022} & نظرية \\
5 & تعلّم معظم المشابك حاصل ضرب أثر محلي في عدد مشترك واحد $\delta$ (المعادلة الثانية في~\eqref{eq:machine}، الفصل~\ref{sec:dofa}) & عصبونات \nt{DOPA} لدى القرد تستجيب لخطأ التنبؤ بالمكافأة؛ وفي شرائح الجسم المخطط يكبّر الدوبامين الشويكة فقط إذا وصل بعد $0.3\text{--}2\,\text{s}$ من الدخل الغلوتاماتي، فالأثر مقيس & \cite{Schultz1997,Yagishita2014} & مُقاس \\
6 & سلك خطأ مستقل إلى كل مخرج لا يوجد إلا في دارة تعلّم الحركات (الفصل~\ref{sec:motorlearning}) & المخيخ: تزامن إشارة الليف المتسلق مع الدخل يُضعف هذا الدخل؛ ولدى القرد تحمل النبضات المعقدة لخلايا بوركنجي اتجاه خطأ الحركة وتستدعي التصحيح & \cite{Ito2001,Herzfeld2018} & مُقاس \\
7 & كل قناة بث حزمة من عدة خطوط (القضية~\ref{prop:bundle}) & في \nt{DOPA}: في المهمة الواحدة تحمل عصبونات مختلفة المكافأة والحركة وسمات المنبّه؛ والخطأ السريع مختلف عن المستوى البطيء للدوبامين. أما في \nt{SERO} و\nt{NORA} و\nt{ACET} فهذا التوزيع أقل دراسة & \cite{Engelhard2019,Hamid2016,Mohebi2019} & مُقاس \\
8 & \nt{SERO}: منظّم المستوى، وبحسب الفرضية الأفق $\tau_\gamma$ (الفصل~\ref{sec:serocomputer}) & التثبيط الذاتي: ناهضات مستقبلاته \foreignlanguage{english}{5-HT$_{1A}$} تثبّط عصبونات السيروتونين في نوى الرفاء، وهذا مُقاس. الأفق مقترح نظريًا؛ وأقوى تجربة أن التنشيط البصري الوراثي لهذه العصبونات لدى الفأر يطيل انتظار المكافأة المؤجلة & \cite{Sprouse1987,Doya2002,Miyazaki2014} & فرضية \\
9 & \nt{NORA}: الكسب $g$ يختار بين البحث والحسم (الفصل~\ref{sec:nora}) & ارتفاع نسبة الإشارة إلى الضجيج: لدى القرد اليقظ يكبت النورأدرينالين في القشرة السمعية النشاطَ الخلفي أكثر مما يكبت الاستجابة لأصوات أبناء جنسه، وهذا مُقاس؛ والكسب الضربي نموذج؛ والدور في الاختيار بين البحث والحسم نظرية & \cite{Foote1975NE,ServanSchreiber1990,AstonJones2005} & فرضية \\
10 & \nt{ACET}: يبدّل بين الكتابة والقراءة (الفصل~\ref{sec:acet}) & الأفعال الثلاثة (إضعاف الوصلات الداخلية، وتقوية الدخل، وتيسير اللدونة) مقيسة؛ وتبديل الأنماط تدعمه بصورة غير مباشرة تجربة السكوبولامين لدى الإنسان & \cite{HasselmoBower1992,Gil1997,Atri2004} & فرضية \\
11 & الصيغة~\eqref{eq:machine} تصف الآلة بكاملها & هي خلاصة نماذج الفصول من~\ref{sec:glutcomputer} إلى~\ref{sec:acet}، وكل جزء منها يستند إلى فصله؛ ولم تُختبر تجريبيًا بوصفها كلًّا & استنتاج في الفصل & فرضية \\
12 & تعمل المقابض منسَّقةً بلا تحكم مشترك: الخطأ، المفاجأة، الكتابة، العودة (الشكل~\ref{fig:timeline}) & لا توجد تجربة: المخطط نوعي. حلقات منفردة: نظرية تقسيم العمل بين \nt{NORA} و\nt{ACET}، وارتباط حدقة العين بسرعة التعلم لدى الإنسان & \cite{YuDayan2005,Nassar2012} & فرضية \\
13 & التعلم بإشارة مشتركة واحدة يبطؤ كلما كثرت العصبونات المؤثرة في النتيجة (القضية~\ref{prop:broadcastcost}) & حساب منحنيات التعلم لشبكة خطية: التعلم باضطراب المخارج أبطأ من الانحدار التدرجي بعدد مرات يساوي عدد المخارج & \cite{Werfel2005} & نظرية \\
14 & كيف تضبط القشرة الدارات المتعددة الطبقات مجهول؛ والمرشحون: رشقات النبضات، والحسابات في فروع العصبون، والتعلم الذاتي بالتنبؤ & الرشقات حاملةً للخطأ نموذج؛ ولم تستطع تكرار استجابة نموذج مفصّل لعصبون قشري إلا شبكة من $5\text{--}8$ طبقات، وهذا نموذج؛ والنبضات الكالسيومية في فروع عصبونات قشرة الإنسان التي تعطي ”أو الحصرية“ مقيسة في الشرائح؛ والتعلم الذاتي مراجعة للأدلة & \cite{Payeur2021,Beniaguev2021,Gidon2020,Keller2018} & فرضية \\
\bottomrule
\end{longtable}}
